% ============================================================================
%  SECCIÓN 3.21: IMPLEMENTACIÓN DE LA NAVEGACIÓN
% ============================================================================

\section{Implementaci\'on de la navegaci\'on}

La navegaci\'on es uno de los aspectos m\'as importantes de una
aplicaci\'on m\'ovil, pues define c\'omo el usuario se desplaza entre las
diferentes pantallas \citep{pressman2010}. En el presente proyecto se
utiliz\'o la librer\'ia \texttt{expo-router}, que implementa un sistema de
navegaci\'on basado en archivos (file-based routing) similar al de Next.js,
donde cada archivo en el directorio \texttt{src/app/} representa una ruta de
la aplicaci\'on.

\subsection{Estructura de navegaci\'on}

La arquitectura de navegaci\'on de la aplicaci\'on se organiza en dos
niveles principales: un \textbf{Stack Navigator} ra\'iz que gestiona las
pantallas de autenticaci\'on y un \textbf{Tab Navigator} que gestiona las
pesta\~nas principales de la aplicaci\'on.

\begin{lstlisting}
src/app/
  _layout.tsx              Stack Navigator ra\'iz
  splash.tsx               Pantalla de bienvenida
  auth/
    login.tsx              Pantalla de inicio de sesi\'on
    register.tsx           Pantalla de registro
  (tabs)/
    _layout.tsx            Tab Navigator
    index.tsx              Pesta\~na Inicio
    search.tsx             Pesta\~na Buscar
    favorites.tsx          Pesta\~na Favoritos
    profile.tsx            Pesta\~na Perfil
  scenario/
    [id].tsx               Detalle del escenario (ruta din\'amica)
\end{lstlisting}

\subsection{Configuraci\'on del Stack Navigator ra\'iz}

El layout ra\'iz (\texttt{\_layout.tsx}) define un \texttt{Stack Navigator}
que gestiona todas las pantallas de la aplicaci\'on. Todas las pantallas se
configuran con \texttt{headerShown: false} para ocultar la barra de
navegaci\'on nativa y utilizar las propias barras dise\~nadas en cada pantalla.

\begin{lstlisting}
export default function RootLayout() {
  return (
    <AuthProvider>
      <AuthGuard>
        <StatusBar style="dark" />
        <Stack screenOptions={{ headerShown: false }}>
          <Stack.Screen name="splash" />
          <Stack.Screen name="(tabs)" />
          <Stack.Screen name="auth/login" />
          <Stack.Screen name="auth/register" />
          <Stack.Screen name="scenario/[id]" />
        </Stack>
      </AuthGuard>
    </AuthProvider>
  );
}
\end{lstlisting}

Los proveedores \texttt{AuthProvider} y el componente \texttt{AuthGuard}
se envuelven alrededor del Stack Navigator para que toda la aplicaci\'on tenga
acceso al contexto de autenticaci\'on y a la protecci\'on de rutas.

\subsection{Configuraci\'on del Tab Navigator}

El Tab Navigator (\texttt{(tabs)/\_layout.tsx}) define las cuatro
pesta\~nas principales de la aplicaci\'on con sus respectivos \'iconos y
t\'itulos. Cada pesta\~na se configura con los colores del tema para
mantener la identidad visual.

\begin{lstlisting}
export default function TabLayout() {
  return (
    <Tabs
      screenOptions={{
        tabBarActiveTintColor: colors.primary,
        tabBarInactiveTintColor: colors.textSecondary,
        headerShown: false,
      }}
    >
      <Tabs.Screen
        name="index"
        options={{
          title: 'Inicio',
          tabBarIcon: ({ color, size }) =>
            <Ionicons name="home-outline" size={size} color={color} />,
        }}
      />
      <Tabs.Screen name="search" options={{ title: 'Buscar', ... }} />
      <Tabs.Screen name="favorites" options={{ title: 'Favoritos', ... }} />
      <Tabs.Screen name="profile" options={{ title: 'Perfil', ... }} />
    </Tabs>
  );
}
\end{lstlisting}

\subsection{Navegaci\'on entre pantallas}

La navegaci\'on entre pantallas se implementa mediante el hook
\texttt{useRouter} de \texttt{expo-router}, que permite ejecutar funciones
de navegaci\'on como \texttt{replace} (reemplazar la pantalla actual) y
\texttt{push} (apilar una nueva pantalla).

\begin{lstlisting}
// Desde Splash hacia Login
const router = useRouter();
const handleComenzar = () => {
  router.replace('/auth/login');
};

// Desde Login hacia Tabs (despu\'es de autenticar)
const onSubmit = async (data) => {
  const { error } = await signIn(data.email, data.password);
  if (!error) router.replace('/(tabs)');
};

// Navegaci\'on con Link (declarativa)
<Link href="/auth/register" asChild>
  <TouchableOpacity>
    <Text style={styles.link}> Registrate</Text>
  </TouchableOpacity>
</Link>
\end{lstlisting}

\subsection{Rutas din\'amicas}

La pantalla de detalle del escenario utiliza una ruta din\'amica
(\texttt{scenario/[id].tsx}) que recibe el identificador del escenario como
par\'ametro de URL. Este patr\'on permite navegar a la detalle de cualquier
escenario con una sola ruta.

\begin{lstlisting}
// Navegaci\'on hacia el detalle de un escenario
router.push('/scenario/123');

// En la pantalla de detalle, lectura del par\'ametro
import { useLocalSearchParams } from 'expo-router';

export default function ScenarioDetailScreen() {
  const { id } = useLocalSearchParams<{ id: string }>();
  // ...
}
\end{lstlisting}

\subsection{Diagrama de navegaci\'on implementado}

\begin{figure}[H]
\centering
  \figuraAPA{fig:navegacion-implementada}{Diagrama de navegaci\'on implementado en la aplicaci\'on \nombreApp{}}{%
  \begin{tikzpicture}[
      nodo/.style={rectangle, draw=black!70, rounded corners=2pt,
                   minimum width=3.0cm, minimum height=0.7cm,
                   align=center, fill=gray!10, font=\scriptsize},
      stack/.style={rectangle, draw=blue!60, rounded corners=3pt,
                    dashed, inner sep=10pt, fill=blue!3},
      tab/.style={rectangle, draw=green!60, rounded corners=3pt,
                  dashed, inner sep=10pt, fill=green!3},
      flecha/.style={-{Stealth[length=2mm]}, thick}
    ]
    % Stack Navigator
    \node[stack, fit={(0,0) (7.5,4)}, label={[font=\small\bfseries]above:Stack Navigator}] (stack) {};
    \node[nodo, fill=yellow!15] (splash) at (3.75,3.0) {Splash};
    \node[nodo] (login) at (1.8,1.5) {Login};
    \node[nodo] (register) at (5.7,1.5) {Registro};
    \node[nodo] (detail) at (5.7,0.0) {Scenario [id]};

    % Tab Navigator
    \node[tab, fit={(10.5,-1.5) (17.5,4)}, label={[font=\small\bfseries]above:Tab Navigator}] (tabs) {};
    \node[nodo] (home) at (11.8,2.8) {Inicio};
    \node[nodo] (search) at (15.5,2.8) {Buscar};
    \node[nodo] (fav) at (11.8,0.8) {Favoritos};
    \node[nodo] (profile) at (15.5,0.8) {Perfil};

    \draw[flecha] (splash) -- (login);
    \draw[flecha] (splash) -- (register);
    \draw[flecha] (login) -- (register) node[midway, below, font=\tiny] {Registrate};
    \draw[flecha] (login) -- (home) node[midway, above, font=\tiny] {Autenticar};
    \draw[flecha] (home) -- (detail);
    \draw[flecha, dashed] (profile) to[bend right=25] node[below, font=\tiny] {signOut} (login);
  \end{tikzpicture}%
  }{Elaboraci\'on propia en base al c\'odigo fuente de la aplicaci\'on.}
\end{figure}
